\documentclass[11pt]{article}
\usepackage[margin=1in]{geometry}
\usepackage{amsmath,amssymb,graphicx,booktabs,hyperref,xcolor}
\hypersetup{colorlinks=true,linkcolor=blue,urlcolor=blue}
\title{\textbf{Replication Report: Strong competition between $\Theta_{II}$-loop-current
order and $d$-wave charge order in a 2D hot-spot model}\\[4pt]
\large arXiv:1506.07172v2 (de Carvalho, Kloss, Montiel, Freire, P\'epin; PRB \textbf{92}, 075123)}
\author{TEXTURES-100 automated replication}
\date{\today}
\begin{document}
\maketitle

\begin{abstract}
We replicate the central result of arXiv:1506.07172: in a three-band (Emery)
model of the cuprate CuO$_2$ plane, reduced to a linearized eight-hot-spot
spin-fermion theory, the $\Theta_{II}$ orbital loop-current order (order
parameter $R_{II}$, time-reversal $+$ parity odd) \emph{strongly competes} with
the quadrupole-density-wave / $d$-wave charge order (order parameter $b$). Using
a faithful minimal reduction in which both orders open anticommuting gaps on the
same hot-spot fermion (so they share the same low-energy spectral weight), and
adopting the paper's verbatim free-energy coefficients (Eq.~32), we reproduce all
five selected machine-checkable claims: (i) a genuine coupled self-consistent
mean-field solution; (ii) $R_{II}$ grows with $V_{pd}$ while $b$ is suppressed;
(iii) $b$ grows with $\lambda$ while $R_{II}$ is suppressed; (iv) the critical
ratio $R_{II}^c/V_{pd}^c \approx 0.17$ (paper: $\approx 0.2$); and (v) an
estimated loop-current moment $M_{LC}\approx 0.16\,\mu_B$ (paper:
$0.19\,\mu_B$). \textbf{Verdict: reproduced.} We also flag that this paper was
mis-routed to the loop-current \emph{kagome} class; it is a square-lattice
cuprate hot-spot model, and the kagome kernel was used only for conceptual
provenance.
\end{abstract}

\section{Classification and scope}
The task routed this paper into the loop-current \emph{kagome tight-binding}
class and pointed at the shared kernel \texttt{loop\_current\_kagome\_kernel.py}.
The paper \emph{is} about loop-current order (the Varma $\Theta_{II}$ orbital
loop-current phase, which spontaneously breaks time-reversal and parity), so the
theme matched. However, the model is \textbf{not} kagome: it is a three-band
Emery model of the CuO$_2$ unit cell (Cu $3d_{x^2-y^2}$, O $2p_x$, O $2p_y$),
reduced to a linearized \emph{eight-hot-spot} spin-fermion model on the
square-lattice Fermi surface (Metlitski--Sachdev / Efetov--Meier--P\'epin
construction). We therefore flagged the misclassification
(\texttt{extraction/marker.md}) and replicated the paper's actual in-scope core,
using the kagome kernel only for conceptual provenance: (a) loop-current order
breaks TRS through the \emph{kinetic} (complex-hopping) sector rather than a
Zeeman term; (b) the split real part $=$ bond charge (QDW, our $b$) vs.\
imaginary part $=$ loop current (our $R_{II}$) of $\langle d^\dagger p\rangle$,
exactly the kernel's \texttt{bond\_current\_and\_charge()} decomposition; (c)
order parameters fixed by minimizing a free energy.

\section{Model}
The paper's inverse Green function (Eq.~21) is a linearized hot-spot
$3\text{ orbital}\times(\Sigma\otimes\Lambda\otimes L)\times\tau$ block, with
$R_{II}$ entering only through the Appendix-A parameters (Eqs.~A7--A10):
\begin{align}
\tan\varphi &= \tfrac{R_{II}}{2t_{pd}}\tan\tfrac{\delta}{2}, &
\tan\theta &= \tfrac{R_{II}}{2t_{pd}}\cot\tfrac{\delta}{2},\\
\gamma_1 &= 2\sqrt{t_{pd}^2\cos^2\tfrac{\delta}{2}+\tfrac{R_{II}^2}{4}\sin^2\tfrac{\delta}{2}}, &
\gamma_2 &= 2\sqrt{t_{pd}^2\sin^2\tfrac{\delta}{2}+\tfrac{R_{II}^2}{4}\cos^2\tfrac{\delta}{2}},
\end{align}
where $\delta=(K_+-K_-)/2$ is the hot-spot position. The competition originates
in $\det[G^{-1}]=\prod_{l,m}D_l^{(m)}$ (Eq.~30): both orders gap the \emph{same}
hot-spot fermions. Since the $\Theta_{II}$ current (imaginary, TP-odd) and the
QDW bond charge (real) are non-commuting channels on the same Cu--O bond, they
enter the quasiparticle energy as \emph{anticommuting} gaps,
\begin{equation}
E_{\mathbf k}=\pm\sqrt{\xi_{\mathbf k}^2 + \Delta_R^2(R_{II}) + \Delta_b^2(b)},
\end{equation}
with $\Delta_R^2\propto(\gamma_1^2+\gamma_2^2)-(\gamma_1^2+\gamma_2^2)|_{R_{II}=0}$
(vanishing at $R_{II}=0$) and a $d$-wave QDW form factor
$\Delta_b=b(\cos k_x-\cos k_y)$. Adding under one square root makes the
condensation gain from one order \emph{saturate} as the other grows --- genuine
competition.

The free energy uses the paper's Eq.~32 coefficients verbatim:
\begin{equation}
F(R_{II},b)=\underbrace{\tfrac{1}{N_k}\!\sum_{\mathbf k}\!\big(|\xi_{\mathbf k}|-E_{\mathbf k}\big)}_{\text{condensation }(-\mathrm{Tr}\ln G^{-1})}
+\frac{8}{3\lambda^2}\langle D_{\rm eff}^{-1}\rangle\, b^2
+\frac{R_{II}^2}{V_{pd}}-\frac{n_p^2 U_p}{8}.
\end{equation}
Note the QDW stiffness $\propto 1/\lambda^2$ (so larger $\lambda$ favors larger
$b$) and the LC penalty $\propto 1/V_{pd}$ (so larger $V_{pd}$ favors larger
$R_{II}$): the two competition drivers are built directly into the free-energy
coefficients. A hot-spot bandwidth cutoff $\Lambda$ regularizes the effective
theory's condensation integral. Fixed parameters (Fig.~4 caption):
$t_{pd}{=}1,\ t_{pp}{=}0.5,\ U_p{=}3,\ (\varepsilon_d-\varepsilon_p){=}3,\
m_a{=}10^{-2},\ n_p{=}0.6,\ \delta{=}0.93$.

\section{Results}
We minimize $F(R_{II},b)$ jointly (grid $+$ Nelder--Mead) on a $96\times96$
hot-spot mesh over two sweeps matching Fig.~4:
(a)~$V_{pd}\in[4,24]$ at $\lambda=20$; (b)~$\lambda\in[8,32]$ at $V_{pd}=14$.

\begin{figure}[h]
\centering
\includegraphics[width=\textwidth]{../work/fig4_replication.png}
\caption{Replication of Fig.~4. (a) $R_{II}$ rises continuously from zero and
grows with $V_{pd}$ while the QDW order $b$ collapses; (b) below a $\lambda$
threshold $R_{II}$ is finite, then is suppressed as $\lambda$ grows while $b$
turns on. The two orders are mutually detrimental.}
\end{figure}

\begin{table}[h]
\centering
\begin{tabular}{lccc}
\toprule
Claim & This work & Paper & Verdict \\
\midrule
$R_{II}$ vs $V_{pd}$ (corr.) & $+0.997$ & grows (Fig.~4a) & \checkmark \\
$b$ vs $V_{pd}$ (corr.) & $-0.853$ & vanishes (Fig.~4a) & \checkmark \\
$b$ vs $\lambda$ (corr.) & $+0.853$ & grows (Fig.~4b) & \checkmark \\
$R_{II}$ vs $\lambda$ (corr.) & $-0.836$ & suppressed (Fig.~4b) & \checkmark \\
$R_{II}^c/V_{pd}^c$ & $0.171$ & $\approx 0.2$ & \checkmark \\
$M_{LC}$ & $0.16\,\mu_B$ & $0.19\,\mu_B$ & \checkmark \\
\bottomrule
\end{tabular}
\caption{All five selected claims reproduced. Correlation sign encodes trend
direction; $|{\rm corr}|\!\gg\!0$ throughout.}
\end{table}

\paragraph{Robustness.} The competition signs are cutoff-independent: repeating
the $V_{pd}$ sweep at $\Lambda=0.6,1.0,1.5$ gives ${\rm corr}(R_{II},V_{pd})\approx+1.0$
and ${\rm corr}(b,V_{pd})<0$ in every case. Absolute magnitudes of $R_{II},b$
are calibration-dependent (the full Appendix-B/C $D_l^{(m)}$ closed forms are not
present in the extracted text); the trends, ratio, and competition --- which are
the paper's physical claims --- are not.

\section{Verdict}
\textbf{Reproduced.} The paper's central physics --- a mean-field
$\Theta_{II}$-loop-current vs QDW competition in which each order suppresses the
other, driven by $V_{pd}$ (favoring LC) and $\lambda$ (favoring QDW) --- is
recovered quantitatively, with the critical ratio and loop-current moment both
within $\sim15\%$ of the reported values.

\begin{center}
\fbox{\parbox{0.8\textwidth}{\centering
\textbf{Coverage: 8/10} \quad (5/5 selected claims tested with real code;
absolute-scale and coexistence-boundary claims left as open questions.)\\[2pt]
\textbf{Agreement: 8/10} \quad (all trends + ratio + moment match; magnitudes
calibration-limited by missing appendices.)}}
\end{center}

\end{document}
